\documentclass[11pt,a4paper]{article}
\usepackage[utf8]{inputenc}
\usepackage{amsmath, amssymb}
\usepackage{booktabs}
\usepackage{geometry}
\geometry{margin=2.5cm}
\usepackage{xurl}
\usepackage[colorlinks=true, linkcolor=black, urlcolor=blue, citecolor=black]{hyperref}

\title{{{title}}}
\author{First Author \and Second Author}
\date{\today}

\begin{document}
\maketitle

\begin{abstract}
State the problem, the approach, and the main result in a few sentences.
\end{abstract}

\section{Introduction}

Introduce the context and the question this article answers. Cite prior work as needed \cite{knuth1984}.

\section{Method}

Describe what was done, in enough detail for someone else to repeat it.

\begin{equation}
    f(x) = \frac{1}{\sigma\sqrt{2\pi}} \, e^{-\frac{(x-\mu)^2}{2\sigma^2}}
    \label{eq:normal}
\end{equation}

Equation~\eqref{eq:normal} is referenced automatically.

\section{Results}

\begin{table}[h]
    \centering
    \begin{tabular}{@{}lcc@{}}
        \toprule
        \textbf{Condition} & \textbf{Mean} & \textbf{SD} \\
        \midrule
        Control   & 0.00 & 1.00 \\
        Treatment & 0.42 & 0.95 \\
        \bottomrule
    \end{tabular}
    \caption{Summary of results.}
    \label{tab:results}
\end{table}

\section{Conclusion}

Summarise the findings and their implications.

\begin{thebibliography}{9}
\bibitem{knuth1984}
    D.~E. Knuth, \emph{The \TeX book}. Addison-Wesley, 1984.
\end{thebibliography}

\end{document}
